\section{Battery-Gated Forward-Only Adaptation}

We consider a controller that can optionally invoke a lightweight forward-only adaptation module before acting. The module is motivated by backprop-free embedding alignment \citep{xiao2026architecture}: it updates a compact latent or embedding transform using only forward statistics, avoiding a full gradient-based inner loop.

At each timestep $t$, the scheduler estimates a robustness gain $\gain$ and an adaptation cost $\adaptcost$. The battery state $\bstate$ determines how aggressively compute cost should be penalized. The scheduling score is
\begin{equation}
s_t = \gain - \lambda(\bstate)\adaptcost,
\end{equation}
where $\lambda(\bstate)$ increases as the remaining battery decreases. The controller adapts only when
\begin{equation}
a_t = \I[s_t > \tau],
\end{equation}
with threshold $\tau$ tuned once for the benchmark.

This formulation is intentionally modest. It is not a new general TTA algorithm. Instead, it isolates a single scheduling decision: whether the controller should spend energy on adaptation right now. The threshold $\tau$ is tuned once on a held-out sweep to maximize a battery-adjusted mission-utility objective. The four policies in the benchmark are then:
\begin{itemize}
\item \textbf{Static}: never adapt.
\item \textbf{Periodic}: adapt on a fixed interval.
\item \textbf{Always-on}: adapt at every step.
\item \textbf{Battery-gated}: adapt only when the score $s_t$ exceeds the threshold.
\end{itemize}

The value of this formulation is that it makes a falsifiable prediction. If adaptation is effectively free, always-on adaptation should dominate. If adaptation is expensive relative to its robustness gain, a gated policy should become preferable, especially in low-battery regimes. The simulator below is designed to test exactly that scheduling boundary.
